Fix one stage and suppress its regularizer in the iteration indices.
The analysis has three steps. An ordinary accelerated energy bounds the
objective gap. The same comparison argument in a different metric bounds
the squared response of the actual projected trajectory. Finally, a
signed-flow inequality converts that response bound into cumulative degree
volume. The second and third steps account for locality.

\subsubsection{One comparison identity, two metrics}

\begin{lemma}[Accelerated comparison]
\label{lem:det-comparison}
In a Hilbert inner product, let $\varphi$ be one-smooth and
$\mu_{\mathrm{c}}$-strongly convex, where $\mu_{\mathrm{c}}=\theta^2$ and
$0<\theta<1$. Write $\chi=1-\theta$ and set
\[
 \vy=\frac{\vx+\theta\vz}{1+\theta},\quad
 \vg=\nabla\varphi(\vy),\quad
 \vq=\chi\vz+\theta\vy-\vg/\theta,\quad
 \vx^+=\chi\vx+\theta\vp.
\]
Here the gradient and norms use the Hilbert metric. If a comparison vector
$\vu$ satisfies
\begin{equation}
 \inner{\vp-\vu}{\vq-\vp}\geq0,
 \label{eq:det-comparison-sector}
\end{equation}
then
\begin{align*}
 &\varphi(\vx^+)-\varphi(\vu)
       +\frac{\mu_{\mathrm{c}}}2\norm{\vp-\vu}^2\\
 &\quad\leq\chi\left(\varphi(\vx)-\varphi(\vu)
       +\frac{\mu_{\mathrm{c}}}2\norm{\vz-\vu}^2\right)
       -\frac{\mu_{\mathrm{c}}\theta(1-\mu_{\mathrm{c}})}2
        \norm{\vz-\vy}^2.
\end{align*}
The vector $\vu$ need not minimize $\varphi$, and the quantity in
parentheses need not be nonnegative.
\end{lemma}
\begin{proof}
The update satisfies $\vx^+=\vy-\vg+\theta(\vp-\vq)$.
Smoothness and expansion of its squared displacement imply
\[
 \varphi(\vx^+)\leq\varphi(\vy)-\tfrac12\norm{\vg}^2
                     +\tfrac{\mu_{\mathrm{c}}}2\norm{\vp-\vq}^2.
\]
The sector inequality gives
$\norm{\vp-\vu}^2\leq\norm{\vq-\vu}^2-\norm{\vp-\vq}^2$.
Add the two inequalities and expand
$\vq-\vu=\chi(\vz-\vu)+\theta(\vy-\vu)-\vg/\theta$.
Strong convexity at $\vx$ and $\vu$, with weights $\chi$ and $\theta$,
respectively, gives
\begin{align*}
 \chi\varphi(\vx)+\theta\varphi(\vu)
 &\geq\varphi(\vy)
    +\inner{\vg}{\chi(\vx-\vy)+\theta(\vu-\vy)}\\
 &\quad+\tfrac{\mu_{\mathrm{c}}}2
       \bigl(\chi\norm{\vx-\vy}^2+\theta\norm{\vu-\vy}^2\bigr).
\end{align*}
The gradient terms cancel because
\[
 \chi(\vx-\vy)+\theta(\vu-\vy)
 =-\theta\{\chi(\vz-\vu)+\theta(\vy-\vu)\}.
\]
Use $\vx-\vy=\theta(\vy-\vz)$ and the convex-combination identity
\[
 \norm{\chi(\vz-\vu)+\theta(\vy-\vu)}^2
 =\chi\norm{\vz-\vu}^2+\theta\norm{\vy-\vu}^2
       -\chi\theta\norm{\vz-\vy}^2.
\]
The remaining decrease is
$-\mu_{\mathrm{c}}\chi\theta(1+\theta)\norm{\vz-\vy}^2/2$.
Since $\chi(1+\theta)=1-\mu_{\mathrm{c}}$, this proves the claim.
\end{proof}

\Needspace{5\baselineskip}
\begin{proposition}[Accelerated objective convergence]
\label{prop:det-stage-rate}
The exact stage recurrence satisfies
\begin{equation}
 E_{k+1}\leq\chi E_k,\qquad
 E_k:=J_r(\vxi_k)-J_r(\vxi^*)
          +\tfrac{\mu_{\mathrm{c}}}2\norm{\vz_k-\vxi^*}_2^2,
 \qquad E_0\leq1.
 \label{eq:det-first-energy}
\end{equation}
Consequently $K\leq1+\theta^{-1}\log(1/\tau)$ steps suffice for
correction gap at most $\tau\in(0,1)$.
\end{proposition}
\begin{proof}
The Hessian of $J_r$ lies between $\alpha\mI$ and $\mI$, and
$\mu_{\mathrm{c}}\leq\alpha$. Euclidean projection onto $\gK_r$ supplies
\cref{eq:det-comparison-sector} with $\vu=\vxi^*$.
Apply \cref{lem:det-comparison} and discard its nonpositive last term.
The correction slack
\begin{equation*}
 \vell_r:=\mQ\vxi^*-\vh+\lambda_r\vomega
         =\mQ\vx_r^*-\vb+\lambda_r\vomega
\end{equation*}
is nonnegative and satisfies $\vell_r^\trans\vxi^*=0$.
Indeed $\vxi^*$ can be positive only where $\vx_r^*$ is positive.
It follows that
\[
 E_0=\tfrac12\norm{\vxi^*}_{\mQ}^2
               +\tfrac{\mu_{\mathrm{c}}}2\norm{\vxi^*}_2^2\leq1,
\]
because $\vomega^\trans\vxi^*\leq1$, $\omega_i\geq1$, and
$\mQ\preceq\mI$. Finally $\chi^K\leq\exp(-\theta K)$.
\end{proof}

\subsubsection{The projection inequality in the PageRank metric}

Euclidean projection onto a box with a mass cap is not generally
nonexpansive in an unrelated matrix metric. We need only the following
inequality for the particular inverse-source comparator $\vt$.

\begin{lemma}[The second projection comparison]
\label{lem:det-sector}
For any $\vq$, put $\vp=\operatorname{Proj}_{\gK_r}(\vq)$ and
$\vn=\vq-\vp$. Then
\begin{equation*}
 \inner{\vp-\vt}{\mQ\vn}
       =(\mQ\vp-\vh)^\trans\vn\geq0.
\end{equation*}
\end{lemma}
\begin{proof}
The Euclidean normal has a decomposition
$\vn=\gamma\vomega+\vn^{\mathrm{up}}-\vn^{\mathrm{low}}$,
where all multipliers are nonnegative and obey complementary slackness.
At a lower face $p_i=0$, Stieltjes signs give
$(\mQ\vp-\vh)_i\leq0$. At an upper face $p_i=4r\omega_i$,
the upper box and the same signs give
$(\mQ\vp)_i\geq4r(\mQ\vomega)_i=4\lambda_r\omega_i\geq h_i$.
Finally, if $\gamma>0$, the mass constraint is tight and
\[
 \vomega^\trans(\mQ\vp-\vh)
 =\alpha\vomega^\trans\vp-m_h=\alpha m_r-m_h=0.
\]
Each normal component has nonnegative pairing with $\mQ\vp-\vh$.
\end{proof}

To bound local work, we need a response estimate that retains the source's
diffuseness. We choose an auxiliary quadratic whose gradient in the $\mQ$
metric is the same update direction as for $J_r$. Comparing it with $\vt$
will make its initial energy proportional to $\lambda_r m_h$. Define
\begin{equation}
 \begin{aligned}
 \mathcal A_r(\vxi)&:=\tfrac12\norm{\mQ\vxi-\vh}_2^2
                  +\alpha\lambda_r\vomega^\trans\vxi,\\
 B_k&:=\mathcal A_r(\vxi_k)-\mathcal A_r(\vt)
                  +\tfrac{\mu_{\mathrm{c}}}2\norm{\vz_k-\vt}_{\mQ}^2.
 \end{aligned}
 \label{eq:det-second-energy}
\end{equation}
Its gradient in the $\mQ$ inner product is
\begin{equation}
 \nabla_{\mQ}\mathcal A_r(\vxi)
 =\mQ^{-1}\{\mQ(\mQ\vxi-\vh)+\alpha\lambda_r\vomega\}
 =\mQ\vxi-\vh+\lambda_r\vomega.
 \label{eq:det-metric-gradient}
\end{equation}
This is exactly the vector already used in \cref{eq:det-iteration}.
The metric Hessian is $\mQ$, so its eigenvalues are in $[\alpha,1]$.

\begin{lemma}[Uniform squared-response bound]
\label{lem:det-response}
For the actual projected trajectory, at every iteration,
\begin{equation*}
 B_{k+1}\leq\chi B_k,\qquad
 \norm{\mQ(\vxi_k-\vxi^*)}_2^2
       \leq18\lambda_r m_h\leq18\alpha^2r.
\end{equation*}
\end{lemma}
\begin{proof}
Apply \cref{lem:det-comparison} in the $\mQ$ metric, using
\cref{lem:det-sector} and \cref{eq:det-metric-gradient}.
The comparator $\vt$ need not minimize $\mathcal A_r$; the comparison
lemma explicitly permits this.
Diffuseness and \cref{lem:det-correction-domain} give
\[
 \norm{\vh}_2^2\leq4\lambda_r m_h,
 \quad \vt^\trans\mQ\vt=\vt^\trans\vh\leq4r m_h,
 \quad \mathcal A_r(\vt)=\lambda_r m_h.
\]
Thus $B_0\leq3\lambda_r m_h$ and
$B_k\leq\chi^k B_0\leq3\lambda_r m_h$, even if $B_0$ is negative.
The linear and kinetic terms in \cref{eq:det-second-energy} are nonnegative,
which yields $\norm{\mQ\vxi_k-\vh}_2^2\leq8\lambda_r m_h$.
Meanwhile
\[
 \vq^*:=\vh-\mQ\vxi^*=\vb-\mQ\vx_r^*
 \quad\hbox{satisfies}\quad
 \vzero\leq\vq^*\leq\lambda_r\vomega,
 \quad \vomega^\trans\vq^*\leq m_h.
\]
Hence $\norm{\vq^*}_2^2\leq\lambda_r m_h$.
The squared triangle inequality gives the factor $2(8+1)=18$.
\end{proof}

\subsubsection{An analytical set with a uniform outside margin}

Let $\gC_r:=\supp(\vx_{r/2}^*)$. This set is used only in the proof.
By \cref{lem:rppr-order-mass},
\begin{equation}
 \vol(\gC_r)\leq2/r,\qquad
 i\notin\gC_r\quad\Longrightarrow\quad
 (\vx_r^*)_i=0,\quad (\vell_r)_i\geq\lambda_r\omega_i/2.
 \label{eq:det-comparison-margin}
\end{equation}
To verify the margin, both optima vanish at such $i$.
The order $\vx_{r/2}^*\geq\vx_r^*$ and the nonpositive off-diagonals
imply $(\mQ\vx_r^*)_i\geq(\mQ\vx_{r/2}^*)_i$.
Subtracting the two regularization loads and applying KKT at $r/2$
proves the claim. Thus even when complementarity is arbitrarily weak on
the optimal boundary, every excursion outside $\gC_r$ pays a known margin.

\begin{theorem}[Cumulative kinetic volume]
\label{thm:det-kinetic-work}
For every horizon $K\geq1$ of an exact stage,
\begin{equation}
 \sum_{k=0}^{K-1}\vol(\supp\vz_{k+1})\leq\frac{76K}{r}.
 \label{eq:det-kinetic-work}
\end{equation}
The sum includes all repeated visits; no confinement assumption is imposed
on the kinetic supports.
\end{theorem}

Thus $K$ accelerated iterations incur $O(K/r)$ total degree volume,
counting every repeated visit. The sparse implementation in
\cref{sec:deterministic-implementation} accounts for the remaining
state and reporting operations.
\begin{proof}
Put $\mK_0=(\mI+\mP)/2$, where $\mP=\mA\mD^{-1}$ is nonnegative
and column stochastic. In mass coordinates set
\[
 \vv_k=\theta\mD^{1/2}\vz_k,\quad
 \vv^*=\theta\mD^{1/2}\vxi^*,\quad
 \beta_0=\frac{1-\alpha}{1+\theta}\leq\chi.
\]
The inequality $\beta_0\leq\chi$ follows from $\theta^2\leq\alpha$.
Let $\vv_k^{\mathrm{raw}}:=\theta\mD^{1/2}\vq_k$ be the proposed mass
before projection. Substitution into the raw update gives the exact identity
\begin{equation*}
 \begin{aligned}
 \vv_k^{\mathrm{raw}}-\vv^*
   &=\beta_0\mK_0(\vv_k-\vv^*)+\vh_k^{\mathrm{f}}-\vr_{\mathrm{f}}^*,\\
 \vh_k^{\mathrm{f}}&=-\frac{\mD^{1/2}(\mQ-\mu_{\mathrm{c}}\mI)
                         (\vxi_k-\vxi^*)}{1+\theta},\qquad
 \vr_{\mathrm{f}}^*=\mD^{1/2}\vell_r.
 \end{aligned}
\end{equation*}
For clarity, the coefficient of $\vv_k$ is
$(\mI-\mD^{1/2}\mQ\mD^{-1/2})/(1+\theta)=\beta_0\mK_0$;
the constant comparison terms cancel by the optimum's stationarity with
slack $\vell_r$.

Select $\gI_k=\{i:v_{k+1,i}>v_i^*\}$ and let
$\ve_k=(\vv_k-\vv^*)_+$. On selected coordinates the projected value is
positive, so lower-face multipliers vanish. Upper-face and mass multipliers
only subtract nonnegative mass. Column stochasticity therefore gives
\begin{equation}
 \norm{\ve_{k+1}}_1+\sum_{i\in \gI_k}r_{{\mathrm{f}},i}^*
 \leq\beta_0\norm{\ve_k}_1+\sum_{i\in \gI_k}h_{k,i}^{\mathrm{f}}.
 \label{eq:det-selected-flow}
\end{equation}
The forcing sum is restricted to the selected coordinates $\gI_k$.
Their degree volume will enter the Cauchy--Schwarz estimate below;
summing all positive forcing over the graph would lose this connection.
Since $\ve_0=\vzero$, summing leaves the nonnegative term
$\norm{\ve_K}_1+(1-\beta_0)\sum_{k=1}^{K-1}\norm{\ve_k}_1$ on the left.
Dropping it bounds cumulative selected slack by cumulative signed forcing.
Every positive kinetic coordinate outside $\gC_r$ is selected because
$\vxi^*$ vanishes there. With
\[
 V_{\mathrm{out}}:=\sum_{k<K}\vol(\supp\vz_{k+1}\setminus\gC_r),
 \qquad H_2:=\sum_{k<K}\norm{\mD^{-1/2}\vh_k^{\mathrm{f}}}_2^2,
\]
we obtain
\begin{equation}
 \tfrac{\lambda_r}2 V_{\mathrm{out}}
 \leq\sum_{k<K}\sum_{i\in \gI_k}h_{k,i}^{\mathrm{f}}
 \leq\sqrt{(V_{\mathrm{out}}+2K/r)H_2}.
 \label{eq:det-selected-cauchy}
\end{equation}
The last inequality is weighted Cauchy--Schwarz: the total selected
degree volume is at most $V_{\mathrm{out}}+2K/r$.
Since $\mQ-\mu_{\mathrm{c}}\mI$ commutes with $\mQ$ and has eigenvalues
between zero and those of $\mQ$, \cref{lem:det-response} implies
$H_2\leq18\alpha^2rK$.
Writing $Y=V_{\mathrm{out}}+2K/r$, \cref{eq:det-selected-cauchy} and
Young's inequality give
\[
 Y\leq\frac{2K}{r}+\frac{2\sqrt{YH_2}}{\lambda_r}
 \leq\frac{2K}{r}+\frac Y2+\frac{2H_2}{\lambda_r^2}.
\]
Thus $Y\leq4K/r+4H_2/\lambda_r^2\leq76K/r$.
The total kinetic volume is at most $Y$.
\end{proof}

An accelerated iteration bound and a small final support would not prove
\cref{eq:det-kinetic-work}. The bound concerns the actual projected
trajectory, including every excursion and revisit. Its usefulness depends
on implementing each iteration in time proportional to kinetic updates
and sparse source records, rather than scanning the historical primal
support. We give that implementation next.
